% Part: lambda-calculus
% Chapter: lambda-definability
% Section: truth-values

\documentclass[../../../include/open-logic-section]{subfiles}

\begin{document}

\olfileid{lam}{ldf}{tvr}
\olsection{رښتينې ارزښتونه او اړيکې}

په خالص لامبډا حساب کښې رښتينې ارزښتونه داسې کوډولے شو:
\begin{align*}
  \fn{true} & \ident \lambd[x][\lambd[y][x]]\\
  \fn{false} & \ident \lambd[x][\lambd[y][y]]
\end{align*}

رښتينې ارزښتونه د \emph{ټاکونکو} په توګه تمثيلېږي؛ يعنې
داسې تابعې دي چې دوه آرګومېنټونه اخلي او يو ئې بېرته
ورکوي۔ رښتينے ارزښت $\fn{true}$ لومړے آرګومېنټ او
$\fn{false}$ دويم آرګومېنټ ټاکي۔ د بېلګې په توګه،
$\fn{true}\, M N$ تل $M$ ته او $\fn{false}\, M N$
تل~$N$ ته راکمېږي۔

\begin{defn}
يوه اړيکه $R \subseteq \Nat^k$ هغه وخت !!{lambda definable}
بولو چې داسې ترم~$R$ وي چې
\begin{align*}
  R\, \num{n_1} \dots \num{n_k} & \bred \fn{true}
  \intertext{هر کله چې $R(n_1, \dots, n_k)$ وي؛ او}
  R\, \num{n_1} \dots \num{n_k} & \bred \fn{false}
\end{align*}
په نورو حالتونو کښې۔
\textbf{د سرچينې سمون (OLLAM-051):} د سرچينې د اړيکې
د ځايونو شمېر د لاندې آرګومېنټونو له شمېر سره نۀ و
برابر؛ دلته د دواړو دپاره يوه نښه کارول شوې ده۔
د اړيکې او د هغې د استازي ترم نوم په اصلي متن کښې يو
شان ليکل شوے؛ له سياقه ئې جلا معنا څرګندېږي۔
\end{defn}

د بېلګې په توګه، اړيکه $\fn{IsZero} = \{0\}$، چې
پر $0$ او يوازې پر $0$ صدق کوي، د لاندې ترم په وسيله
!!{lambda definable} ده:
\[
  \fn{IsZero} \ident \lambd[n][n (\lambd[x][\fn{false}])\, \fn{true}].
\]
دا څنګه کار کوي؟ د چرچ عددنښې تکراروونکي دي، يعنې
خپل لومړے آرګومېنټ پر دويم $n$ ځله تطبيقوي۔ دلته
لومړنی ارزښت $\fn{true}$ ږدو او د هر تکرار په ګام کښې
د مخکني ګام له پايلې پرته $\fn{false}$ ورکوو۔ دا ګام
پر لومړني ارزښت $n$ ځله تطبيقېږي؛ پايله هغه وخت او
يوازې هغه وخت $\fn{true}$ وي چې ګام هېڅ تطبيق نۀ شي،
يعنې کله چې $n = 0$ وي۔

د رښتينو ارزښتونو د دې تمثيل پر بنسټ نورې رښتينې تابعې
هم تعريفولے شو۔ دلته دوه دي: د نفي او عطف تمثيلونه:
\begin{align*}
  \fn{Not} & \ident \lambd[x][x\, \fn{false}\, \fn{true}]\\
  \fn{And} & \ident \lambd[x][\lambd[y][xy \,\fn{false}]]
\end{align*}
تابع «$\fn{Not}$» يو آرګومېنټ اخلي؛ که آرګومېنټ
$\fn{false}$ وي، $\fn{true}$ ورکوي، او که آرګومېنټ
هم~$\fn{true}$ وي، $\fn{false}$ ورکوي۔ تابع
«$\fn{And}$» دوه رښتينې ارزښتونه اخلي او بايد هله،
او يوازې هله، $\fn{true}$ ورکړي چې دواړه آرګومېنټونه
هم~$\fn{true}$ وي۔ رښتينې ارزښتونه لکه پورته د ټاکونکو
په توګه تمثيلېږي؛ نو که $x$ رښتينے ارزښت وي او پر دوو
آرګومېنټونو تطبيق شي، پايله لومړے آرګومېنټ وي که
$x$، $\fn{true}$ وي، او کنه دويم آرګومېنټ۔ اوس
$\fn{And}$ خپل دوه آرګومېنټونه $x$ او $y$ اخلي او
بيا $y$ او $\fn{false}$ خپل لومړي آرګومېنټ~$x$ ته
ورکوي۔ که $x$ رښتينے ارزښت وي، پايله هغه وخت~$y$
وي چې $x$، $\fn{true}$ وي؛ او هغه وخت $\fn{false}$
وي چې $x$، $\fn{false}$ وي۔ همدا مو غوښتل۔

پام وکړئ دلته فرض کوو چې $\fn{And}$ ته يوازې رښتينې
ارزښتونه د آرګومېنټونو په توګه ورکول کېږي۔ که نور
ترمونه ورکړل شي، پايله ئې، يعنې عادي بڼه که موجوده
وي، ښايي رښتينے ارزښت نۀ وي۔

\begin{prob}
  د شاملې بېلېوالې او ځانګړې بېلېوالې د رښتينو تابعو
  تمثيلوونکې تابعې $\fn{Or}$ او $\fn{Xor}$ تعريف کړئ؛
  رښتينې ارزښتونه د $\lambd$-ترمونو په توګه کوډ کړئ۔
\end{prob}

\end{document}
